% ============================================================
\section{Personalisation Strategy}
% ============================================================

\subsection{Overall Approach: Rule-Based with Online Adaptation}

Smart Inbox uses a rule-based personalisation strategy rather than a
narrative one. The reason is mainly that email is not really a
narrative medium. There is no story structure or teaching path for the
system to move through. It is just a stream of separate messages
arriving over time. Because of that, it made more sense to score each
email against a fixed set of features and then update the per-user
weights as more behaviour data is collected. This also keeps the logic
visible. The system can explain why a message was ranked in a certain
way, which would be harder to do if that decision sat inside a
narrative layer.

\subsection{Content Adaptation}

The main content adaptation technique is
\textbf{conditional inclusion}~\cite{Brusilovsky2007ch18}. The content
is all there in the system, but it is not always shown in the same
form. What the user sees depends on both the current mode and the view
they are in. In Busy Mode, the full email body is hidden and replaced
with a generated summary and an action-item list. In Normal Mode, the
full body is shown. The interface adds another layer to this. The list
view only shows sender, subject, and snippet, while the detail view
reveals more of the email. So the same message is exposed differently
depending on context.

This works well here because it stops the interface from spending
space on content the user is less likely to need at that moment. The
trade-off is that something can be missed. A user in Busy Mode may
read the summary and miss detail or tone from the original message. To
reduce that risk, the mode toggle stays visible in the navigation bar
so the user can switch back easily.

The system does not use stretchtext. That technique makes more sense
in structured documents where small parts can be expanded from
meaningful hotwords. Email is usually much less structured, so it is
not a very natural fit here. Fragment variants and page variants are
also not really used in the strict sense, although Busy Mode and
Normal Mode do behave like a broad page-level variant.

\subsection{Navigation Adaptation}

The system also changes navigation by using
\textbf{priority-zoned direct guidance}~\cite{Brusilovsky2007nav}
instead of simple chronological sorting. The four tabs guide the user
towards \textit{Now} first, but they do not lock the rest of the inbox
away. In that sense the design is similar to the module's ``Next''
link idea, except the thing being guided is not a sequence of pages
but a set of email buckets.

Lower-priority mail is \textbf{de-emphasised} rather than hidden. The
\textit{Later} tab stays visible and its contents can still be read in
full. This avoids the main weakness of link hiding, where the system
draws too hard a line between relevant and irrelevant and can leave
the user unaware of what was pushed out. De-emphasis is a better fit
here because it reduces distraction without removing awareness.

The explanation shown underneath each email in the detail view works
as \textbf{link annotation}~\cite{Brusilovsky2007nav}. It adds context
about why the message was placed in a particular bucket, so the user
can judge the system's reasoning before deciding whether to trust it
or give feedback. This fits the taxonomy well because the structure
stays stable while the system adds extra guidance on top. In Smart
Inbox this appears through the score breakdown display.

\begin{figure}[H]
\centering
    \includegraphics[height=0.4\textheight, width=\textwidth, keepaspectratio]{images/scoring pipeline.png}
    \caption{Scoring pipeline (left) and bucket threshold gauge (right).}
    \label{fig::uml_diagram-1}
\end{figure}

% ============================================================
\section{System Architecture}
% ============================================================

Smart Inbox follows a three-tier client-server structure, but most of
the adaptive behaviour sits in the backend. Using the module's
architecture, the main pieces are the User Model Repository, the Rules
and Content Repository, and the Adaptive Application itself.
Table~\ref{tab:archmap} shows how the Smart Inbox implementation maps
onto those parts.

\begin{table}[H]
\centering\small
\caption{Mapping of Smart Inbox components to the module's adaptive architecture framework.}
\label{tab:archmap}
\setlength{\tabcolsep}{4pt}
\begin{tabularx}{\linewidth}{L{3.8cm}L{7.6cm}}
\toprule
\textbf{Module component} & \textbf{Smart Inbox implementation} \\
\midrule
\rowcolor{lightrow}
User Model Repository & \texttt{UserPreference}, \texttt{SenderProfile}, \texttt{ThreadProfile}, \texttt{InteractionEvent} tables in SQLite, managed by \texttt{user\_model.py} \\
Rules and Content Repository & Feature weight defaults and scoring rules in \texttt{ranking.py}; email bodies and metadata in the \texttt{emails} table \\
\rowcolor{lightrow}
Adaptive Application & FastAPI backend orchestrating \texttt{adaptation.py} (update logic), \texttt{summary.py} (content generation), and \texttt{routes.py} (API); React frontend rendering personalised views \\
\bottomrule
\end{tabularx}
\end{table}

The frontend uses React 18, TypeScript, and Vite. It renders the
inbox, switches between modes, and sends user actions through
\texttt{api.ts}.

The backend uses FastAPI and Python. That is where most of the
adaptive logic sits. \texttt{auth.py} handles OAuth and token
refresh. \texttt{gmail\_sync.py} performs incremental sync through
Gmail history IDs. \texttt{ranking.py} extracts features and scores
emails. \texttt{adaptation.py} updates the model from events and
feedback. \texttt{user\_model.py} manages weights and profile
helpers. \texttt{summary.py} produces either LLM-based or heuristic
summaries. \texttt{sync.py} provides the fallback path for demo mode
when Gmail authentication is not being used.

The database uses SQLAlchemy over SQLite. It stores eight tables:
\texttt{users}, \texttt{emails}, \texttt{user\_preferences},
\texttt{sender\_profiles}, \texttt{thread\_profiles},
\texttt{interaction\_events}, \texttt{oauth\_states}, and
\texttt{oauth\_tokens}. Together these hold both the message data and
the user-model data that the adaptation logic depends on.

\begin{figure}[H]
    \centering
    \includegraphics[height=0.4\textheight, width=\textwidth, keepaspectratio]{images/Architecture Diagram adaptive.png}
    \caption{System architecture. The LLM service is disabled by default, so the standard configuration keeps processing local.}
    \label{fig:arch}
\end{figure}

\textbf{Data flow:} the data flow is fairly direct. On sync, emails
are fetched from the Gmail API, parsed, scored by \texttt{ranking.py}
using the user's current weight vector, and stored with a bucket
assignment. User interactions are written into
\texttt{InteractionEvent}. Those events are then rolled up into the
sender, thread, and preference profiles. On the next sync, the
updated profile values feed back into the scoring step. That closes
the adaptation loop shown in Figure~\ref{fig:loop}.

\begin{figure}[H]
    \centering
    \includegraphics[height=0.3\textheight, width=\textwidth, keepaspectratio]{images/Adaptation loop.png}
\caption{Adaptation loop and scrutability chain. The orange arc shows model updates feeding back into scoring on the next sync.}
\label{fig:loop}
\end{figure}